%=============================================================
\chapter{System Design}
\label{ch:arch}
%=============================================================

% ─── 章節破題 ───────────────────────────────────────────────
第~\ref{ch:motivation}~章的分析揭示了兩個獨立但相互影響的問題：
其一，beam search 中大多數 hop 是可消除的無效工作（O1、O2）；
其二，收斂偵測的可靠性依賴於搜尋初期的 top-$k$ 覆蓋品質（O3）。
這兩個問題共同決定了 PaceANN 的設計空間：
\textbf{只改動 ET 不改動快取，終止信號不可靠；
只改動快取不改動 ET，I/O 節省有上限。}
兩者必須協同設計，才能構成一個自洽的系統。

本章按以下順序展開：
首先建立快取選擇策略的形式基礎，並說明為何靜態覆蓋不足、必須採用動態快取（§\ref{sec:topology}）；
其次描述並發感知動態快取的設計（§\ref{sec:cache}），
再正式定義收斂感知早停（ET）條件及其可靠性前提（§\ref{sec:et}），
接著論證快取與 ET 的協同必要性（§\ref{sec:codesign}），
最後說明 server 模式下的並發架構（§\ref{sec:server}）。

%-------------------------------------------------------------
\section{快取節點選擇的拓撲基礎}
\label{sec:topology}
%-------------------------------------------------------------

設計快取策略的第一個問題是：
在 DRAM 預算有限的條件下，\emph{哪些節點的鄰居列表應被優先快取？}
這個問題等價於估計各節點的存取頻率，並按頻率遞減排序選取。

\paragraph{存取頻率的拓撲決定性}

在 DiskANN 的 beam search 中，
每次查詢都從固定入口節點（medoid $m^*$）出發，
沿 Vamana 圖的有向邊貪婪地向 query 方向逼近。
設 $d_G(v, m^*)$ 為節點 $v$ 到 medoid 在圖上的最短路徑長度（跳數）。

\begin{proposition}[存取頻率單調遞減性]
\label{prop:monotone}
設查詢向量服從資料集上的均勻分布，
令 $f(v)$ 為節點 $v$ 的期望存取次數。
則 $f(v)$ 是 $d_G(v, m^*)$ 的嚴格單調遞減函數：
\[
  d_G(u, m^*) < d_G(v, m^*) \implies f(u) > f(v)
\]
\end{proposition}

\begin{proof}[直觀論證]
對任意查詢 $q$，beam search 的路徑必然起始於 $m^*$（$d_G = 0$）。
$d_G = 1$ 的節點（$m^*$ 的直接鄰居）會在第一個 hop 全部被納入候選，
因此被訪問的概率接近 1。
隨著 $d_G$ 增加，到達節點 $v$ 的路徑逐漸分叉——
不同查詢在第 $d_G$ 跳時已沿不同方向延伸，
只有路徑恰好經過 $v$ 的查詢才會訪問它，
故 $f(v)$ 隨 $d_G$ 的增大呈指數衰減。
\end{proof}

\paragraph{命題的設計推論}

命題~\ref{prop:monotone} 有兩個直接推論：

\begin{enumerate}
  \item \textbf{medoid 鄰域是全域最高頻的區域，且可在服務任何查詢之前被識別。}
    最靠近 medoid 的節點其存取頻率完全由圖拓撲決定，與查詢分布無關，
    因此系統啟動時（index 載入後）即可精確識別，無需任何查詢歷史。

  \item \textbf{BFS-from-medoid 是最佳的\emph{冷啟動種子}策略。}
    以廣度優先搜尋（BFS）從 $m^*$ 展開、按層序選取節點，
    恰好等價於按 $d_G(v, m^*)$ 由小到大排序——即按期望訪問頻率由高到低選取。
    因此在快取尚無查詢歷史的冷啟動階段，以 BFS-from-medoid 預填快取，
    能立即保護最高頻的入口鄰域（隨機採樣則會讓大量低頻節點佔據預算，故次優；
    這也解釋了 Starling 的隨機採樣 navigation graph 為何命中率偏低）。
\end{enumerate}

\paragraph{為何種子不足以取代動態快取}

然而，medoid 鄰域的高頻性\emph{並不}意味著靜態覆蓋就足夠。
每個查詢真正耗時的節點集中在\textbf{該查詢的收斂區}——
即 beam search 逼近 query 時、答案附近才被展開的節點；
這些節點是\emph{查詢專屬}的，且會隨查詢分布（例如 RAG 的熱門主題）持續漂移，
無法被任何在系統啟動時固定的集合覆蓋（第~\ref{ch:motivation}~章 O3）。
換言之，BFS 種子蓋得住\emph{入口熱區}，卻蓋不住\emph{各查詢的收斂熱區}。
這正是 PaceANN 不採用永久靜態快取、而以\textbf{並發感知動態快取}（§\ref{sec:cache}）
在查詢期間依實際存取分布持續演化的原因；BFS-from-medoid 僅作為其冷啟動種子。
對多 medoid 索引（$|M| = K$），種子採\textbf{分片均等預算}：
總預算平均分配為 $N/K$ 給每個 medoid 各自 BFS，避免部分入口鄰域在冷啟動時完全未覆蓋。

%-------------------------------------------------------------
\section{並發感知動態快取（BNC）}
\label{sec:cache}
%-------------------------------------------------------------

如 §\ref{sec:topology} 所述，靜態覆蓋只能保護入口熱區，
無法涵蓋各查詢專屬、且隨分布漂移的收斂熱區。
因此 PaceANN 的快取以\textbf{需求驅動的動態快取}為主體
（BoundedNeighborCache，BNC），在固定記憶體預算內於查詢期間持續演化，
自動反映實際工作負載的存取模式；
系統啟動時以 §\ref{sec:topology} 的 BFS-from-medoid 種子預填，消弭冷啟動初期的空窗，
其後即交由 BNC 的替換策略託管，平滑演化至當前 workload 的分布。

\paragraph{設計約束}

BNC 的設計面臨四個並行存在的約束：

\begin{enumerate}
  \item \textbf{高並發讀取}：$T = 16$ 個 worker thread 同時查找，
    全域鎖（global lock）會形成嚴重序列化瓶頸。
  \item \textbf{低插入頻率}：每個節點首次觸及 SSD 後插入一次，
    之後只有被驅逐時才需再次插入，插入路徑可以承擔更高的成本。
  \item \textbf{固定記憶體上限}：需精確控制 DRAM 使用量，
    不能因為工作負載波動而超出預算。
  \item \textbf{近似 LRU 語義}：淘汰最不常被近期訪問的節點，
    但不需要精確 LRU 的 $O(1)$ doubly-linked list 維護。
\end{enumerate}

\paragraph{分片架構：256 個獨立 CLOCK 環}

BNC 以 \texttt{node\_id \% 256} 將節點分配至 256 個 shard，
每個 shard 獨立維護一個 CLOCK 替換環與一個 \texttt{shared\_mutex}：

\begin{itemize}
  \item \textbf{查找（高頻）}：取 \texttt{shared\_lock}（允許並發讀者），
    雜湊表查找 + atomic reference bit 設定，整體約 10--20\,ns。
    256 個 shard 使得鎖競爭降低至 $1/256$，
    16 個 thread 平均每個 shard 的並發讀者只有 $\sim\!0.06$ 個。

  \item \textbf{插入與驅逐（低頻）}：取 \texttt{unique\_lock}（互斥），
    掃描 CLOCK hand 找到 reference bit 為 false 的 slot，
    原地替換，成本為 $O(\text{shard capacity} / \text{hit rate})$，
    不影響同 shard 的查找路徑（因已持有 unique lock）。
\end{itemize}

相較於精確 LRU，CLOCK 省去了每次查找時更新 doubly-linked list 的
exclusive lock 操作，在並發讀取遠多於插入的場景下
（正是 server mode 的典型情況），
吞吐量與尾延遲均優於 LRU。

\paragraph{座標共存（Coordinate Co-Storage）}

標準 DiskANN 的搜尋流程分為兩個階段：
（1）traversal：以 PQ 距離引導的 beam search，只需鄰居 IDs；
（2）re-rank：取 top-$k$ 候選的完整原始向量，計算精確距離。
Re-rank 階段需要從 SSD 讀取 $K$ 個節點的完整向量，
即使 traversal 已全部命中快取，也仍有不可避免的 SSD 存取。

BNC 在快取項目中\textbf{共同儲存節點的原始座標}（raw uint8 coordinates），
使命中 BNC 的節點在 re-rank 時可直接從 DRAM 讀取精確向量，
完全消除 re-rank 的 SSD I/O：

\begin{center}
\texttt{struct Entry \{ uint32\_t* neighbors;\ uint8\_t* coords;\ bool ref; \}}
\end{center}

此設計使得「BNC 命中的節點貢獻零 SSD I/O」在整個搜尋流程中成立，
而非僅在 traversal 階段成立。
代價是每個 BNC entry 需要額外儲存 $d_\text{raw}$ bytes 的座標（例如 128 bytes for uint8 SIFT），
在記憶體預算配置時需一併納入計算。

\paragraph{需求驅動填充與冷啟動}

除了啟動時的 BFS 種子外，BNC 採用\textbf{需求驅動（demand-fill）}策略：
節點的鄰居列表在首次從 SSD 讀取後才被填入，而非批量預載。
如此 BNC 自動反映實際工作負載的存取模式，並隨查詢分布的變化持續適應。
冷啟動初期的入口空窗由 BFS 種子（§\ref{sec:topology}）填補，
其餘的收斂熱區則在服務過程中透過 demand-fill 逐步累積、由 CLOCK 替換策略託管；
因此系統\emph{不依賴}預先的代表性查詢暖機，即可在真實流量下平滑收斂至穩態命中率。

%-------------------------------------------------------------
\section{PQ 引導的自適應早停}
\label{sec:et}
%-------------------------------------------------------------

早停的核心邏輯在 §\ref{sec:mot-insight} 已引入；
本節給出正式定義，並說明關鍵設計參數的含義。

\paragraph{停止條件：跨尺度精確門檻}

在 beam search 每次展開前，取 retset 中 PQ 下界最小的未展開候選 $v_\text{best\text{-}unexp}$，
將其 PQ 下界 $d_\text{PQ}(q, v_\text{best\text{-}unexp})$ 與\emph{已讀回全精度向量}算得的、
目前第 $2k$ 名候選的\textbf{精確}距離 $d_\text{exact}(q, v_{2k})$ 相比：

\begin{equation}
  d_\text{PQ}(q,\, v_\text{best\text{-}unexp}) \;>\; \theta \cdot d_\text{exact}(q,\, v_{2k})
  \label{eq:et-gate}
\end{equation}

若成立，代表 frontier 中最有希望候選的\emph{樂觀}下界都已劣於現有第 $2k$ 名的\emph{實際}距離，
繼續展開幾乎不可能改變 top-$k$，即可終止並返回。
此設計刻意採「跨尺度」比較——待展開端用便宜的 PQ 下界、已收斂端用精確距離——
並以第 $2k$ 名（而非第 $k$ 名）為門檻，為 recall 預留保守邊界。
$d_\text{exact}$ 使用 traversal 過程中已讀回的全精度向量（含 BNC co-storage 命中者），
因此此判斷\emph{不額外觸發任何 SSD I/O}。

\paragraph{$\theta$ 的語義}

$\theta$ 控制停止條件的保守程度：

\begin{itemize}
  \item $\theta = 1.0$：frontier 中最佳候選的 PQ 距離等於 top-$k$ 邊界時即停止，
    對 PQ 近似完全信任，較為激進。
  \item $\theta \to \infty$：永不停止，等同原始 DiskANN。
  \item 典型設定（$\theta \approx 1.1$\textasciitilde$1.3$）：
    為 PQ 量化誤差預留容忍帶，
    在保守程度與 hop 節省之間取得平衡。
\end{itemize}

$\theta$ 的最優值依賴資料集的 PQ 壓縮誤差分布，
不同 dimension 與 segment 數量的 PQ 設定下，
$\theta$ 的最優點不同。
參數敏感性分析在第~\ref{ch:method}~章中呈現。

\paragraph{Patience 與最小 hop 數安全網}

單一 hop 的 PQ 下界可能因量化噪聲短暫越過門檻，
而搜尋最初期（hop 1\textasciitilde 3）top-$k$ 尚未有足夠候選、$d_\text{exact}(q, v_{2k})$ 亦不穩定，
兩者都可能造成假陽性終止。系統以兩道安全網防範：

\begin{equation}
  \text{Stop iff } \big(\text{式}~\eqref{eq:et-gate}\ \text{連續成立 } P \text{ 個 hop}\big)
  \ \text{AND}\ t \geq t_\text{min}
  \label{eq:et-guard}
\end{equation}

其一是 \textbf{patience 計數}：唯有式~\eqref{eq:et-gate} 連續在 $P$ 個 hop 都成立才真正終止
（典型 $P \approx 6$），過濾單跳噪聲造成的偽觸發。
其二是\textbf{最小 hop 數} $t_\text{min}$：在前 $t_\text{min}$ 個 hop 內一律不允許 ET，
確保 top-$k$ 累積到足夠穩定後才開始判斷。
兩者皆不依賴任何靜態快取層——收斂端的精確距離由 traversal 已讀回、
並由動態快取 co-storage 保存的全精度向量提供。

%-------------------------------------------------------------
\section{快取與早停的協同必要性}
\label{sec:codesign}
%-------------------------------------------------------------

本節論證快取與 ET 的協同不是「兩個獨立優化的疊加」，
而是一個設計中的相互依賴關係：
兩者缺少任何一個，另一個都會出現功能性的退化。

\paragraph{情況一：有快取無 ET}

若只部署動態快取（BNC），
beam search 仍按固定 $L$ 執行完整的 hop 序列。
此時所有 hop 的 SSD 讀取被快取命中替換為 DRAM 存取（~10\,ns），
traversal 成本大幅降低，但\textbf{hop 數本身未減少}。

在 §\ref{sec:mot-diminishing} 的分析中，
$L = 150$ 時有 70.7\% 的 hop 在收斂後才執行。
有快取時，這些無效 hop 仍會消耗 DRAM 頻寬、BNC shard 鎖，
以及 CPU 的 PQ 距離計算資源。
在高並發下，DRAM 頻寬競爭與鎖競爭會成為新的瓶頸，
原本被 SSD I/O 掩蓋的競爭現象在快取命中率提高後反而顯現。

\paragraph{情況二：有 ET 無快取（冷啟動）}

若只部署 ET，而沒有快取保護前期 hop，
系統在搜尋剛起步時（hop 1\textasciitilde 3）的 top-$k$ 覆蓋不足，
$d_\text{PQ}(q, v_{k\text{th}})$ 不穩定，
ET 的比值 $r_t$ 容易出現假陽性，迫使 $\theta$ 設定得非常保守。

在無快取的冷啟動條件下，
即使 $\theta = 1.2$，假陽性率仍可能使 recall 損失超出可接受範圍，
導致 $\theta$ 必須調大到 $\sim\!1.4$ 以上才能維持目標 recall。
但此時 ET 的觸發率極低，hop 節省幅度縮水，
相較於 baseline DiskANN 的改善微乎其微。

\paragraph{協同效應：快取讓 ET 可靠，ET 讓快取的價值充分實現}

兩者協同時：

\begin{enumerate}
  \item 動態快取（含 BFS 種子與 demand-fill 累積的熱點）使前期 top-$k$ 以 DRAM 速度快速填充，
    並透過 co-storage 直接提供收斂端所需的\emph{精確}距離 $d_\text{exact}$。
  \item top-$k$ 穩定後（通常在 hop 5\textasciitilde 6），$d_\text{exact}(q, v_{2k})$ 可靠，
    ET 的跨尺度門檻得以在較小的 $\theta$ 下穩定觸發。
  \item ET 在收斂點停止，消除後續的無效 hop——
    這些 hop 即使在快取命中的情況下也是純粹的 DRAM 與 CPU 浪費。
  \item 快取在收斂區持續命中，使 ET 之前的有效 hop 更快完成，進一步縮短整體延遲。
\end{enumerate}

在此協同下，一個 query 的服務路徑可形式化為：

\begin{equation}
  \text{Latency} = \underbrace{h^* \cdot c_\text{DRAM}}_{\text{有效 hop（動態快取）}}
                 + \underbrace{0}_{\text{無效 hop（ET 阻止）}}
                 + \underbrace{K \cdot c_\text{SSD\_rerank}}_{\text{re-rank（有 coord-co-storage 時 = 0）}}
  \label{eq:latency-model}
\end{equation}

其中 $h^*$ 是 oracle 停止 hop，$c_\text{DRAM}$ 是每次快取命中的成本。
這個模型代表系統在理想條件下的延遲下界，
實際延遲高於此下界的原因主要是 BNC 命中率不為 100\%（仍有部分 SSD 讀取），
以及 shard 鎖競爭的開銷。

%-------------------------------------------------------------
\section{Server 模式並發架構}
\label{sec:server}
%-------------------------------------------------------------

\paragraph{執行緒模型}

PaceANN server 採用固定執行緒池設計：
$T$ 個 worker thread 各自持有一個 TCP socket 監聽，
以及一個 thread-local scratch space（用於 retset、visited bitmap 等 per-query 狀態），
避免跨 query 的狀態污染。
每個連線被接受後，由空閒 worker 執行完整的搜尋流程後關閉連線，
無需跨 thread 的工作竊取或 continuation 機制。

\paragraph{在制查詢計數（Active Search Counter）}

BNC 的 shard 鎖在高並發下的競爭程度取決於同時執行中的 query 數量。
server 維護一個原子計數器 \texttt{\_active\_searches}（\texttt{std::atomic<int32\_t>}），
在 \texttt{cached\_beam\_search} 進入時遞增、返回後遞減，
用於為動態 I/O 並發控制提供負載信號（見下段）。

\paragraph{I/O 並發控制（io\_concurrency）}

在高並發時（$C \gg T$，連線排隊），
所有 $T$ 個 thread 同時執行 beam search，
每個 hop 最多提交 $W$ 個 SSD 讀取，
使 $D_\text{eff}$ 在快取未命中時可達 $T \times W$。
為控制 NVMe 佇列深度，
系統提供 \texttt{--io\_concurrency N} 選項：

\begin{equation}
  W_\text{eff} = \min\!\left(W,\; \max\!\left(1,\; \left\lfloor \frac{N}{A} \right\rfloor\right)\right)
  \label{eq:io-concurrency}
\end{equation}

其中 $A = \texttt{\_active\_searches}$，
$N$ 為目標最大並發 SSD 讀取數。
當 $A = T = 16$、$N = 64$ 時，$W_\text{eff} = 4$，
每 hop 的 SSD 請求數下降 50\%，
以小幅 recall 代價換取 NVMe 佇列的穩定。
在 BNC 充分暖機的情況下，$W_\text{eff}$ 的縮減對實際 SSD 請求量的影響較小
（因多數 hop 命中快取，frontier 中真正需要 SSD 的節點本就不多），
主要作用是作為快取冷啟動期間的保護機制。

\paragraph{二進位 TCP 通訊協定}

client 透過 TCP 連線傳送固定格式的請求封包，
包含查詢向量（float 或 uint8）、$K$、$L$、$\theta$ 與 hop\_budget；
server 返回 top-$K$ 的 node IDs 與精確距離。
所有計時（mean、P50、P99、P999 latency）在 server 端從\textbf{接收完整請求後}開始計算，
不含 TCP 建連與傳輸時間，使延遲量測反映搜尋本身的計算成本。

%-------------------------------------------------------------
\section{設計摘要}
\label{sec:design-summary}
%-------------------------------------------------------------

表~\ref{tab:design-summary} 彙整 PaceANN 各元件的設計選擇與理論依據：

\begin{table}[h]
\centering
\caption{PaceANN 元件設計摘要}
\label{tab:design-summary}
\begin{tabular}{p{3.2cm} p{3.8cm} p{5.5cm}}
\toprule
元件 & 設計選擇 & 理論依據 \\
\midrule
動態快取冷啟動 & BFS-from-medoid 種子（分片均等預算） & 存取頻率單調遞減性（命題~\ref{prop:monotone}）；僅作冷啟動 \\
快取替換策略 & CLOCK（256 shard） & 高並發讀取下 CLOCK > LRU（省去 unique-lock 更新）\\
快取座標共存 & 鄰居 ID 與全精度向量共存 & 消除 re-rank 的 SSD I/O，並供 ET 精確距離 \\
ET 停止條件 & 跨尺度精確門檻 $d_\text{PQ} > \theta\, d_\text{exact}(2k)$ & 零額外 I/O 的收斂信號 \\
ET 安全網 & patience $P$ + 最小 hop $t_\text{min}$ & 防止單跳噪聲與 top-$k$ 未穩定的假陽性 \\
I/O 並發控制 & $W_\text{eff} = N / A$ & 動態平衡 NVMe 佇列深度與搜尋寬度 \\
\bottomrule
\end{tabular}
\end{table}
